\section{Method}
\label{sec:method}

\subsection{Device structure}

\begin{figure*}[t]
\centering\includegraphics[width=\textwidth]{panel01.png}
\caption{Structure and normalization. (a) The simulated cross-section to scale: a
\SI{5}{\nano\metre} silicon nanosheet body gated on both faces through
\SI{1}{\nano\metre} \ce{SiO2} and \SI{7}{\nano\metre} HZO, with a
\SI{100}{\nano\metre} gate and \SI{15}{\nano\metre} source and drain overlaps. The
marked point is where polarization and electric field are probed in every transient
reported here. (b) The mapping from the two-dimensional slab to the nanosheet. Because
the slab is gated on both faces, a slab of depth $A$ carries gate width $2A$; matching the
nanosheet perimeter $2(W+T_\mathrm{si})=\SI{90}{\nano\metre}$ fixes
$A=\SI{0.045}{\micro\metre}$. Runs executed at \num{0.071}; the correction is an exact
post-multiplier, verified against a re-simulated node to \num{1.8e-16}. Computed values:
\fignote{A1}.}
\label{fig:panel1}
\end{figure*}

The simulated device is a silicon nanosheet of thickness $T_\mathrm{si}=\SI{5}{\nano\metre}$
and width $W=\SI{40}{\nano\metre}$, gated through a \SI{1}{\nano\metre} interfacial
\ce{SiO2} layer and a \SI{7}{\nano\metre} \ce{Hf}\ce{Zr}\ce{O_x} (HZO) ferroelectric,
with a \SI{5}{\nano\metre} TiN gate metal of work function \SI{4.35}{\electronvolt}.
The gate is \SI{100}{\nano\metre} long and overlaps the source and drain extensions by
\SI{15}{\nano\metre} on each side. The body is doped
$N_\mathrm{sub}=\SI{1e16}{\per\cubic\centi\metre}$ and the source and drain
$N_\mathrm{sd}=\SI{5e19}{\per\cubic\centi\metre}$. Figure~\ref{fig:panel1}(a) shows the stack to scale,
together with the position $(0,\,\SI{7}{\nano\metre})$ at which polarization and
electric field are probed in every transient reported here.

This geometry is the endpoint of a coordinate-descent search over read bias,
$T_\mathrm{ox}$, $T_\mathrm{si}$, $N_\mathrm{sub}$, $T_\mathrm{fe}$ and gate length,
each turn evaluated on the retained on/off window at the sub-coercive operating point.
The trajectory is in Fig.~\ref{fig:panel8}(a) and the individual sweeps in
Figs.~\ref{fig:panel8}(b)--(c).

\subsection{Simulation setup}

Structures were built in Sentaurus Structure Editor and solved in Sentaurus Device.
The ferroelectric is the Preisach model, whose parameters were calibrated as described
in Section~\ref{sec:cal}. Transport uses the Philips unified mobility model with normal-field
degradation. All transient drive is applied as instantaneous hops between short
fixed-bias holds rather than finely stepped ramps.

That last choice is not cosmetic. Under relative-error timestep control, a
non-converging deck does not fail. It \emph{crawls}: the timestep walks down to
whatever minimum is specified and never recovers, so the run continues to produce
output indefinitely. Three decks in this work hit that behaviour and produced log
files of \SI{71}{\mega\byte}, \SI{143}{\mega\byte} and \SI{344}{\mega\byte}, the last
covering a few percent of its intended ramp in 320\,000 steps. The fixes are
structural: instantaneous bias transitions, a minimum timestep floored in proportion
to the maximum, and a wall-clock limit per job so that a remaining crawl is recorded
as a failure instead of consuming a shared license quietly.

\subsection{Calibration}
\label{sec:cal}

The deck was calibrated against published measurements of a gate-all-around
metal--ferroelectric--metal--insulator--semiconductor FeFET. The calibration searched
over remanent and saturation polarization, coercive field, permittivity, interface trap
density, fixed interface charge and the Preisach kinetic prefactor, minimizing the
disagreement between the simulated and measured transfer characteristics on both
polarization states. The locked parameter set is
$P_r=\SI{32}{\micro\coulomb\per\square\centi\metre}$,
$P_s=\SI{40}{\micro\coulomb\per\square\centi\metre}$,
$F_c=\SI{1.4}{\mega\volt\per\centi\metre}$, $\varepsilon=33$,
$D_\mathrm{it}=\SI{4e12}{\per\square\centi\metre}$ and
$Q_f=\SI{7e12}{\per\square\centi\metre}$.

Figure~\ref{fig:panel2} shows the result three ways: the overlay itself, the same curves with the
window and the current level divided out, and the residual.

\paragraph*{What is calibrated, and what is not}
The overlay applies a single free width scalar, fitted so that the erased branch's
saturation current matches the measurement's. That scalar absorbs any constant error in
normalization, and it is large: the simulated saturation current is
\SI{2.150e-4}{\ampere\per\micro\metre} against the measurement's
\SI{1.523e-6}{\ampere\per\micro\metre}, a factor of 141. Consequently this work claims a
calibrated memory window, threshold voltage, subthreshold slope, on/off ratio and
turn-on shape. It does \emph{not} claim that the absolute current density is
experimentally validated, and no conclusion in this paper depends on it: every
device-level result is a ratio, a voltage, a polarization or a field, and the network
results normalize the conductance ladder by its own maximum.

\subsection{From a two-dimensional cross-section to a nanosheet}
\label{sec:norm}

Sentaurus Device solves the two-dimensional cross-section per unit depth in $z$. An
area factor $A$ declares that the slab is $A$ micrometres deep, so a contact current
in the output file is $I=A\,I_\mathrm{2D}$; it is a pure post-multiplier on contact
currents and charges and never enters the Poisson or continuity solve.

The simulated slab is gated on both faces, so a slab of depth $A$ carries a total gate
width $2A$. Matching the real nanosheet's gate perimeter,
\begin{equation}
  W_\mathrm{eff} = \mathrm{TESW} = 2\,(W + T_\mathrm{si})
                 = 2\,(40 + 5)\ \mathrm{nm} = \SI{90}{\nano\metre},
  \label{eq:tesw}
\end{equation}
gives $2A = \SI{0.090}{\micro\metre}$ and therefore $A=\SI{0.045}{\micro\metre}$.
All currents in this paper are reported per micrometre of gate perimeter under
(\ref{eq:tesw}). Figure~\ref{fig:panel1}(b) draws the mapping.

The runs were executed at $A=\SI{0.071}{\micro\metre}$, a value inherited from an
earlier campaign and never re-derived from the geometry. Because $A$ is a post-multiplier,
the correction is an exact rescale by $0.045/0.071=0.6338$ rather than a re-simulation.
This was verified rather than asserted: one node was re-run at $A=\SI{0.045}{\micro\metre}$
and divided by the same node at \num{0.071}, giving \num{0.633802816901} against an
expected \num{0.633802816901}: a relative error of \num{1.8e-16}, which is
double-precision epsilon. Every ratio, voltage, subthreshold slope, polarization and
field quantity is unchanged by the correction; every absolute current, conductance,
resistance, charge and energy scales by \num{0.6338}.

\paragraph*{Limitation: this is a two-dimensional model}
What the check above establishes is that the arithmetic of the mapping is exact. It does
not establish that a flat double-gated slab behaves like a wrapped gate. A real
gate-all-around nanosheet has corners, and the electrostatics at those corners, along with
the edge fields along the sheet width, are not represented in a cross-section, so
corner-enhanced inversion, corner-dependent field concentration in the ferroelectric,
and any width-direction non-uniformity of switching are outside what this model can
show. No three-dimensional gate-all-around simulation was performed. The effect is
expected to matter most for absolute current and least for the ratio and voltage
quantities this paper reports, but that expectation is an argument, not a measurement,
and it is stated here so that it is not mistaken for one.

One second-order consequence is worth naming. The area factor was held at
\num{0.071} across the entire $T_\mathrm{si}$ sweep, while the true perimeter under
(\ref{eq:tesw}) runs from \SI{90}{\nano\metre} to \SI{110}{\nano\metre} across that
sweep, a \SI{22}{\percent} spread the constant factor does not capture. Ratio metrics
cancel it; absolute currents from geometry sweeps do not, and are normalized by each
device's own perimeter where they are used.

\subsection{Write and read protocols}

Three protocols appear in this paper and they are not interchangeable.

\emph{Retained-state characterization} (panels 3--4, 8(d), 9) writes with a
\SI{\pm2}{\volt}, \SI{5}{\micro\second} pulse, then reads with a single continuous
\SI{500}{\nano\second} gate ramp. The read is analysed on the \emph{conduction} current,
the sum of the electron and hole contact currents, so that the ramp's displacement term
through the \SI{15}{\nano\metre} gate--drain overlap is removed exactly rather than
argued to be negligible.

\emph{Analog programming} (panels 5--7, 10) applies trains of identical sub-coercive
pulses at \SI{2.0}{\volt} for \SI{100}{\nano\second}, with a read between pulses at
$V_G=0$ and $V_{DS}=\SI{0.05}{\volt}$. Depression uses \SI{-2.0}{\volt} pulses of at
least \SI{1}{\micro\second}: at \SI{100}{\nano\second} the erase pulse sees only about
\SI{26}{\percent} of the ferroelectric field that a programming pulse of the same
amplitude sees, because the depolarization field opposes erase, so a short erase pulse
does not switch the device and the resulting apparent asymmetry is a protocol artefact.

\emph{Material characterization} (panel 3(a)) uses a standalone TiN/HZO/TiN capacitor
with no semiconductor, swept triangularly. The capacitor is simulated with the
calibrated material value of the depolarization time constant rather than the
application value; on a quasi-static loop the application value bleeds remanent
polarization away between coercive crossings and understates the material.

The depolarization time constant in the application parameter set is
$\tau_P=\SI{10}{\micro\second}$. Every retained value in this paper is read after at
least $5\tau_P$, and each figure states the delay it used.

\subsection{Extraction rules}
\label{sec:rules}

Each of the following rules exists because its absence produced a wrong number in this
work. They are stated so that the numbers below can be reproduced, and because several
of them are the reason a result reported here differs from an earlier one.

\begin{enumerate}
\item \textbf{A ratio may only be formed from two measurements taken in the same run.}
  Three previously reported numbers divided a reading from one run by a reading from
  another and were wrong as a result. The read contrast after nine potentiation pulses,
  for instance, is $18.4\times$ against that run's own baseline and $30.6\times$ against
  a different run's.

\item \textbf{A retained value must be read after the state has settled, and the delay
  must be stated.} Reading immediately after a write measures the depolarization
  transient. Doing so produced three separate fabricated failures in this work: a
  retention collapse, a read-disturb collapse of tens of percent, and an endurance
  collapse, each of which looked like responsible bad news. Measured after settling,
  the same runs give a flat plateau, \SI{0.0000}{\percent} read disturb, and no
  degradation at all.

\item \textbf{A threshold voltage means nothing without its extraction criterion.} The
  same calibration data gives a memory window of \SI{1.296}{\volt} at a constant current
  of \SI{1e-8}{\ampere} per nanosheet and \SI{1.218}{\volt} at
  \SI{1e-7}{\ampere\per\micro\metre}, roughly a decade deeper in subthreshold. Both are
  correct. Every window and threshold quoted here carries its criterion.

\item \textbf{Subthreshold slope is fitted, not read off a pair of points.} On a
  \SI{50}{\milli\volt} grid a single adjacent pair can span a decade, so one point sets
  the answer; the programmed branch's ambipolar minimum reads as
  \SI{45}{\milli\volt\per\text{dec}} that way, which would be an apparently
  sub-thermal slope produced by a branch crossover. Slopes here are the steepest
  straight-line fit over a window spanning at least two decades and starting at least
  one decade above the branch minimum.

\item \textbf{A simulation that completes without an error is not a result.} Two runs in
  this work completed cleanly with essentially zero field everywhere: a parameter set
  with a \SI{100}{\micro\second} depolarization time constant made a
  \SI{500}{\micro\second} write hold five time constants long, so the model relaxed back
  inside the same hold and the probe read \SI{2.3267e-7}{\coulomb\per\square\centi\metre}
  erased against \SI{2.3238e-7}{\coulomb\per\square\centi\metre} programmed, a
  \SI{0.1}{\percent} difference against a remanent polarization of
  \SI{3.2e-5}{\coulomb\per\square\centi\metre}. Nothing from those runs is used here.

\item \textbf{Counts are counted.} The synapse count of the deployed network was
  previously stated as \num{10700}; enumerating the model's parameters gives
  \num{261440} weights, because the recurrent layer is 160 neurons rather than 100 and
  both delayed layers store ten independently programmed taps per connection.

\item \textbf{A figure may not display a number it did not compute.} Every figure in this
  paper is produced by a single generator that writes, alongside each image, the table
  of exactly the values plotted; annotations are computed from those arrays. One
  previously published figure carried a memory window typed in as text. It happened to
  be correct, and it could never have caught an error.
\end{enumerate}

\paragraph*{Two on/off ratios, and which is which}
This device has two legitimate current-ratio figures of merit and they differ by more
than two orders of magnitude. The \emph{retained memory window} is the ratio between the
fully erased and fully programmed states read at $V_G=0$ after settling: $5410\times$.
The \emph{read contrast} is the ratio between the rest state and the state after nine
sub-coercive potentiation pulses: $18.4\times$. The first describes a binary memory
operation, the second an integrate-and-fire event. They are never quoted adjacently in
this paper without that distinction.

\subsection{Energy measurement}

Energy per programming pulse is taken as $V_\mathrm{pgm}\,\Delta Q_G$, with $\Delta Q_G$
read from the solver's own gate-charge trace across each write window. The apparently
more direct route (integrating $V_G I_G$ over the pulse) is not usable at the
available time resolution. The gate current during a write is a switching spike followed
by a flat conduction floor of about \SI{23}{\pico\ampere}, and the transient output is
sampled every \SI{50}{\nano\second}, so the spike is one sample wide and its apparent
width is set by the sampler rather than by the device. Integrating on that grid returns
roughly \SI{20}{\femto\joule} per pulse, three orders of magnitude above the
charge-based value, and it is measuring the sampling interval.

\subsection{Network deployment}

The classifier is a spiking recurrent network trained on single-lead electrocardiogram
beats, with three input channels, a recurrent layer of 100 leaky integrate-and-fire and
60 adaptive leaky integrate-and-fire neurons, and four output classes. Input and
recurrent connections carry a ten-tap delay line, giving \num{261440} weights and
\num{522880} devices in a differential pair arrangement.

Weights are deployed by post-training quantization onto the measured conductance ladder.
Each weight is realized as $(G^+ - G^-)/g_\mathrm{max}$, so a global scale on the
measured current cancels exactly; this is why correcting the area factor, which changes
every absolute current by \num{0.6338}, reproduces the previously locked accuracy to four
decimal places. The checkpoint is fixed throughout: no retraining, quantization-aware or
otherwise, was performed after the device export was corrected.

Cycle-to-cycle noise is injected per conductance branch rather than per weight. Applying
it to the weight understates the noise for weights near zero, where $G^+\approx G^-$ and
the difference is small but each branch's jitter is not.
